% Options for packages loaded elsewhere
% Options for packages loaded elsewhere
\PassOptionsToPackage{unicode}{hyperref}
\PassOptionsToPackage{hyphens}{url}
\PassOptionsToPackage{dvipsnames,svgnames,x11names}{xcolor}
%
\documentclass[
  letterpaper,
  DIV=11,
  numbers=noendperiod]{scrartcl}
\usepackage{xcolor}
\usepackage{amsmath,amssymb}
\setcounter{secnumdepth}{-\maxdimen} % remove section numbering
\usepackage{iftex}
\ifPDFTeX
  \usepackage[T1]{fontenc}
  \usepackage[utf8]{inputenc}
  \usepackage{textcomp} % provide euro and other symbols
\else % if luatex or xetex
  \usepackage{unicode-math} % this also loads fontspec
  \defaultfontfeatures{Scale=MatchLowercase}
  \defaultfontfeatures[\rmfamily]{Ligatures=TeX,Scale=1}
\fi
\usepackage{lmodern}
\ifPDFTeX\else
  % xetex/luatex font selection
\fi
% Use upquote if available, for straight quotes in verbatim environments
\IfFileExists{upquote.sty}{\usepackage{upquote}}{}
\IfFileExists{microtype.sty}{% use microtype if available
  \usepackage[]{microtype}
  \UseMicrotypeSet[protrusion]{basicmath} % disable protrusion for tt fonts
}{}
\makeatletter
\@ifundefined{KOMAClassName}{% if non-KOMA class
  \IfFileExists{parskip.sty}{%
    \usepackage{parskip}
  }{% else
    \setlength{\parindent}{0pt}
    \setlength{\parskip}{6pt plus 2pt minus 1pt}}
}{% if KOMA class
  \KOMAoptions{parskip=half}}
\makeatother
% Make \paragraph and \subparagraph free-standing
\makeatletter
\ifx\paragraph\undefined\else
  \let\oldparagraph\paragraph
  \renewcommand{\paragraph}{
    \@ifstar
      \xxxParagraphStar
      \xxxParagraphNoStar
  }
  \newcommand{\xxxParagraphStar}[1]{\oldparagraph*{#1}\mbox{}}
  \newcommand{\xxxParagraphNoStar}[1]{\oldparagraph{#1}\mbox{}}
\fi
\ifx\subparagraph\undefined\else
  \let\oldsubparagraph\subparagraph
  \renewcommand{\subparagraph}{
    \@ifstar
      \xxxSubParagraphStar
      \xxxSubParagraphNoStar
  }
  \newcommand{\xxxSubParagraphStar}[1]{\oldsubparagraph*{#1}\mbox{}}
  \newcommand{\xxxSubParagraphNoStar}[1]{\oldsubparagraph{#1}\mbox{}}
\fi
\makeatother

\usepackage{color}
\usepackage{fancyvrb}
\newcommand{\VerbBar}{|}
\newcommand{\VERB}{\Verb[commandchars=\\\{\}]}
\DefineVerbatimEnvironment{Highlighting}{Verbatim}{commandchars=\\\{\}}
% Add ',fontsize=\small' for more characters per line
\usepackage{framed}
\definecolor{shadecolor}{RGB}{241,243,245}
\newenvironment{Shaded}{\begin{snugshade}}{\end{snugshade}}
\newcommand{\AlertTok}[1]{\textcolor[rgb]{0.68,0.00,0.00}{#1}}
\newcommand{\AnnotationTok}[1]{\textcolor[rgb]{0.37,0.37,0.37}{#1}}
\newcommand{\AttributeTok}[1]{\textcolor[rgb]{0.40,0.45,0.13}{#1}}
\newcommand{\BaseNTok}[1]{\textcolor[rgb]{0.68,0.00,0.00}{#1}}
\newcommand{\BuiltInTok}[1]{\textcolor[rgb]{0.00,0.23,0.31}{#1}}
\newcommand{\CharTok}[1]{\textcolor[rgb]{0.13,0.47,0.30}{#1}}
\newcommand{\CommentTok}[1]{\textcolor[rgb]{0.37,0.37,0.37}{#1}}
\newcommand{\CommentVarTok}[1]{\textcolor[rgb]{0.37,0.37,0.37}{\textit{#1}}}
\newcommand{\ConstantTok}[1]{\textcolor[rgb]{0.56,0.35,0.01}{#1}}
\newcommand{\ControlFlowTok}[1]{\textcolor[rgb]{0.00,0.23,0.31}{\textbf{#1}}}
\newcommand{\DataTypeTok}[1]{\textcolor[rgb]{0.68,0.00,0.00}{#1}}
\newcommand{\DecValTok}[1]{\textcolor[rgb]{0.68,0.00,0.00}{#1}}
\newcommand{\DocumentationTok}[1]{\textcolor[rgb]{0.37,0.37,0.37}{\textit{#1}}}
\newcommand{\ErrorTok}[1]{\textcolor[rgb]{0.68,0.00,0.00}{#1}}
\newcommand{\ExtensionTok}[1]{\textcolor[rgb]{0.00,0.23,0.31}{#1}}
\newcommand{\FloatTok}[1]{\textcolor[rgb]{0.68,0.00,0.00}{#1}}
\newcommand{\FunctionTok}[1]{\textcolor[rgb]{0.28,0.35,0.67}{#1}}
\newcommand{\ImportTok}[1]{\textcolor[rgb]{0.00,0.46,0.62}{#1}}
\newcommand{\InformationTok}[1]{\textcolor[rgb]{0.37,0.37,0.37}{#1}}
\newcommand{\KeywordTok}[1]{\textcolor[rgb]{0.00,0.23,0.31}{\textbf{#1}}}
\newcommand{\NormalTok}[1]{\textcolor[rgb]{0.00,0.23,0.31}{#1}}
\newcommand{\OperatorTok}[1]{\textcolor[rgb]{0.37,0.37,0.37}{#1}}
\newcommand{\OtherTok}[1]{\textcolor[rgb]{0.00,0.23,0.31}{#1}}
\newcommand{\PreprocessorTok}[1]{\textcolor[rgb]{0.68,0.00,0.00}{#1}}
\newcommand{\RegionMarkerTok}[1]{\textcolor[rgb]{0.00,0.23,0.31}{#1}}
\newcommand{\SpecialCharTok}[1]{\textcolor[rgb]{0.37,0.37,0.37}{#1}}
\newcommand{\SpecialStringTok}[1]{\textcolor[rgb]{0.13,0.47,0.30}{#1}}
\newcommand{\StringTok}[1]{\textcolor[rgb]{0.13,0.47,0.30}{#1}}
\newcommand{\VariableTok}[1]{\textcolor[rgb]{0.07,0.07,0.07}{#1}}
\newcommand{\VerbatimStringTok}[1]{\textcolor[rgb]{0.13,0.47,0.30}{#1}}
\newcommand{\WarningTok}[1]{\textcolor[rgb]{0.37,0.37,0.37}{\textit{#1}}}

\usepackage{longtable,booktabs,array}
\usepackage{calc} % for calculating minipage widths
% Correct order of tables after \paragraph or \subparagraph
\usepackage{etoolbox}
\makeatletter
\patchcmd\longtable{\par}{\if@noskipsec\mbox{}\fi\par}{}{}
\makeatother
% Allow footnotes in longtable head/foot
\IfFileExists{footnotehyper.sty}{\usepackage{footnotehyper}}{\usepackage{footnote}}
\makesavenoteenv{longtable}
\usepackage{graphicx}
\makeatletter
\newsavebox\pandoc@box
\newcommand*\pandocbounded[1]{% scales image to fit in text height/width
  \sbox\pandoc@box{#1}%
  \Gscale@div\@tempa{\textheight}{\dimexpr\ht\pandoc@box+\dp\pandoc@box\relax}%
  \Gscale@div\@tempb{\linewidth}{\wd\pandoc@box}%
  \ifdim\@tempb\p@<\@tempa\p@\let\@tempa\@tempb\fi% select the smaller of both
  \ifdim\@tempa\p@<\p@\scalebox{\@tempa}{\usebox\pandoc@box}%
  \else\usebox{\pandoc@box}%
  \fi%
}
% Set default figure placement to htbp
\def\fps@figure{htbp}
\makeatother





\setlength{\emergencystretch}{3em} % prevent overfull lines

\providecommand{\tightlist}{%
  \setlength{\itemsep}{0pt}\setlength{\parskip}{0pt}}



 


\usepackage{graphicx}
\usepackage{etoolbox}
\AtBeginEnvironment{tabular}{\resizebox{\textwidth}{!}}
\KOMAoption{captions}{tableheading}
\makeatletter
\@ifpackageloaded{caption}{}{\usepackage{caption}}
\AtBeginDocument{%
\ifdefined\contentsname
  \renewcommand*\contentsname{Table of contents}
\else
  \newcommand\contentsname{Table of contents}
\fi
\ifdefined\listfigurename
  \renewcommand*\listfigurename{List of Figures}
\else
  \newcommand\listfigurename{List of Figures}
\fi
\ifdefined\listtablename
  \renewcommand*\listtablename{List of Tables}
\else
  \newcommand\listtablename{List of Tables}
\fi
\ifdefined\figurename
  \renewcommand*\figurename{Figure}
\else
  \newcommand\figurename{Figure}
\fi
\ifdefined\tablename
  \renewcommand*\tablename{Table}
\else
  \newcommand\tablename{Table}
\fi
}
\@ifpackageloaded{float}{}{\usepackage{float}}
\floatstyle{ruled}
\@ifundefined{c@chapter}{\newfloat{codelisting}{h}{lop}}{\newfloat{codelisting}{h}{lop}[chapter]}
\floatname{codelisting}{Listing}
\newcommand*\listoflistings{\listof{codelisting}{List of Listings}}
\makeatother
\makeatletter
\makeatother
\makeatletter
\@ifpackageloaded{caption}{}{\usepackage{caption}}
\@ifpackageloaded{subcaption}{}{\usepackage{subcaption}}
\makeatother
\usepackage{bookmark}
\IfFileExists{xurl.sty}{\usepackage{xurl}}{} % add URL line breaks if available
\urlstyle{same}
\hypersetup{
  pdftitle={Project title},
  colorlinks=true,
  linkcolor={blue},
  filecolor={Maroon},
  citecolor={Blue},
  urlcolor={Blue},
  pdfcreator={LaTeX via pandoc}}


\title{Project title}
\author{Zhalae Daneshvari, Nick Johnson, Jolene Le}
\date{}
\begin{document}
\maketitle


\section{Data 1}\label{data-1}

\subsection{Problem or question}\label{problem-or-question}

\begin{itemize}
\item
  \textbf{Question:} Can we predict what type of adverse reaction a
  patient will experience when taking a cardiovascular drug, based on
  patient demographics and drug information?
\item
  \textbf{Objective:} This project will build a machine learning
  classifier using data from the FDA Adverse Event Reporting System
  (FAERS) to predict the type of reported reaction for common
  cardiovascular medications.
\item
  \textbf{Importance:} This topic is important because cardiovascular
  drugs---such as propranolol, lisinopril, losartan, enalapril, and
  atenolol, are among the most prescribed worldwide and are associated
  with a wide range of side effects that vary in severity and frequency
  across patient groups. Understanding which types of reactions are more
  common for specific drugs and demographics can improve post-market
  surveillance, support safer prescribing practices, and contribute to
  public health insights
\item
  \textbf{Variables:} The dataset includes variables such as drug name
  (categorical), patient age and sex (quantitative and categorical), and
  the reported reaction (categorical target variable), along with
  additional metadata like report year and seriousness indicators.
\item
  \textbf{Deliverables:} The major deliverables will include a cleaned
  and structured dataset derived from FAERS data, and a trained
  classification model that predicts reaction categories.
\end{itemize}

\subsection{Introduction and data}\label{introduction-and-data}

The dataset I compiled includes five widely prescribed cardiovascular
drugs:

\begin{itemize}
\item
  Propranolol (beta blocker)
\item
  Lisinopril, Losartan, Enalapril (ACE inhibitors/ARBs)
\item
  Atenolol (beta blocker)
\end{itemize}

The target variable will be the reaction column, which lists the main
adverse event reported for each case. The model will learn patterns
between demographic and drug-related predictors and the corresponding
reaction type.

\textbf{Data Source}

\begin{itemize}
\item
  \textbf{Dataset:} Custom FAERS subset for five cardiovascular drugs
  (Propranolol, Lisinopril, Losartan, Enalapril, Atenolol)
\item
  \textbf{Source:} openFDA Drug Adverse Event API *
\item
  \textbf{Collector:} U.S. Food and Drug Administration (FDA)
\item
  \textbf{Collection Method: R}eports submitted by healthcare
  professionals, manufacturers, and patients describing post-market
  adverse drug events.
\end{itemize}

\textbf{Data Description}

Each observation in the dataset represents a single adverse event report
submitted to the FDA for one of five cardiovascular drugs: propranolol,
lisinopril, losartan, enalapril, or atenolol. Each record contains basic
patient demographics (age, sex), the specific drug associated with the
report, and one or more reported reactions, which serve as the outcome
variable for this study. The reactions vary widely, including both mild
symptoms (e.g., dizziness, cough, fatigue) and more severe physiological
effects (e.g., hypotension, bradycardia). Additional metadata such as
the seriousness of the report and the year it was filed provide temporal
and contextual insight. Overall, the dataset captures a diverse sample
of post-market adverse event experiences across several classes of
cardiovascular medications, enabling exploration of how reaction types
differ by drug and patient profile.

\textbf{Ethical Considerations}

\begin{itemize}
\item
  \textbf{Privacy:} Data are fully de-identified and publicly available.
\item
  \textbf{Bias:} Reports are voluntarily submitted and may overrepresent
  severe or unusual events.
\item
  \textbf{Interpretation:} Model predictions reflect associations, not
  causal effects.
\end{itemize}

\subsection{Glimpse of data}\label{glimpse-of-data}

\begin{Shaded}
\begin{Highlighting}[]
\ImportTok{import}\NormalTok{ pandas }\ImportTok{as}\NormalTok{ pd}
\ImportTok{import}\NormalTok{ numpy }\ImportTok{as}\NormalTok{ np}
\ImportTok{import}\NormalTok{ ast}

\CommentTok{\# ============================}
\CommentTok{\# 1️⃣ Load dataset}
\CommentTok{\# ============================}
\NormalTok{df }\OperatorTok{=}\NormalTok{ pd.read\_csv(}\StringTok{"data{-}raw/faers\_sample.csv"}\NormalTok{)}

\CommentTok{\# Sample smaller portion for speed}
\NormalTok{df\_sample }\OperatorTok{=}\NormalTok{ df.sample(n}\OperatorTok{=}\BuiltInTok{min}\NormalTok{(}\DecValTok{2000}\NormalTok{, }\BuiltInTok{len}\NormalTok{(df)), random\_state}\OperatorTok{=}\DecValTok{42}\NormalTok{)}


\CommentTok{\# ============================}
\CommentTok{\# 2️⃣ Helper functions to flatten nested JSON{-}like columns}
\CommentTok{\# ============================}
\KeywordTok{def}\NormalTok{ extract\_reaction(text):}
    \CommentTok{"""Extract first reaction term from nested FAERS \textquotesingle{}patient.reaction\textquotesingle{} field."""}
    \ControlFlowTok{try}\NormalTok{:}
\NormalTok{        data }\OperatorTok{=}\NormalTok{ ast.literal\_eval(}\BuiltInTok{str}\NormalTok{(text))}
        \ControlFlowTok{if} \BuiltInTok{isinstance}\NormalTok{(data, }\BuiltInTok{list}\NormalTok{) }\KeywordTok{and} \BuiltInTok{len}\NormalTok{(data) }\OperatorTok{\textgreater{}} \DecValTok{0}\NormalTok{:}
\NormalTok{            val }\OperatorTok{=}\NormalTok{ data[}\DecValTok{0}\NormalTok{]}
            \ControlFlowTok{if} \BuiltInTok{isinstance}\NormalTok{(val, }\BuiltInTok{dict}\NormalTok{) }\KeywordTok{and} \StringTok{"reactionmeddrapt"} \KeywordTok{in}\NormalTok{ val:}
                \ControlFlowTok{return}\NormalTok{ val[}\StringTok{"reactionmeddrapt"}\NormalTok{]}
        \ControlFlowTok{return} \VariableTok{None}
    \ControlFlowTok{except} \PreprocessorTok{Exception}\NormalTok{:}
        \ControlFlowTok{return} \VariableTok{None}


\KeywordTok{def}\NormalTok{ extract\_drug(text):}
    \CommentTok{"""Extract first medicinal product name from nested FAERS \textquotesingle{}patient.drug\textquotesingle{} field."""}
    \ControlFlowTok{try}\NormalTok{:}
\NormalTok{        data }\OperatorTok{=}\NormalTok{ ast.literal\_eval(}\BuiltInTok{str}\NormalTok{(text))}
        \ControlFlowTok{if} \BuiltInTok{isinstance}\NormalTok{(data, }\BuiltInTok{list}\NormalTok{) }\KeywordTok{and} \BuiltInTok{len}\NormalTok{(data) }\OperatorTok{\textgreater{}} \DecValTok{0}\NormalTok{:}
\NormalTok{            val }\OperatorTok{=}\NormalTok{ data[}\DecValTok{0}\NormalTok{]}
            \ControlFlowTok{if} \BuiltInTok{isinstance}\NormalTok{(val, }\BuiltInTok{dict}\NormalTok{) }\KeywordTok{and} \StringTok{"medicinalproduct"} \KeywordTok{in}\NormalTok{ val:}
                \ControlFlowTok{return}\NormalTok{ val[}\StringTok{"medicinalproduct"}\NormalTok{]}
        \ControlFlowTok{return} \VariableTok{None}
    \ControlFlowTok{except} \PreprocessorTok{Exception}\NormalTok{:}
        \ControlFlowTok{return} \VariableTok{None}


\CommentTok{\# Apply flattening functions}
\ControlFlowTok{if} \StringTok{"patient.reaction"} \KeywordTok{in}\NormalTok{ df\_sample.columns:}
\NormalTok{    df\_sample[}\StringTok{"reaction\_clean"}\NormalTok{] }\OperatorTok{=}\NormalTok{ df\_sample[}\StringTok{"patient.reaction"}\NormalTok{].}\BuiltInTok{apply}\NormalTok{(extract\_reaction)}

\ControlFlowTok{if} \StringTok{"patient.drug"} \KeywordTok{in}\NormalTok{ df\_sample.columns:}
\NormalTok{    df\_sample[}\StringTok{"drug\_clean"}\NormalTok{] }\OperatorTok{=}\NormalTok{ df\_sample[}\StringTok{"patient.drug"}\NormalTok{].}\BuiltInTok{apply}\NormalTok{(extract\_drug)}

\CommentTok{\# ============================}
\CommentTok{\# 3️⃣ General dataset overview}
\CommentTok{\# ============================}
\BuiltInTok{print}\NormalTok{(}\StringTok{"📊 Dataset Overview"}\NormalTok{)}
\BuiltInTok{print}\NormalTok{(}\StringTok{"="} \OperatorTok{*} \DecValTok{40}\NormalTok{)}
\BuiltInTok{print}\NormalTok{(}\StringTok{"Shape of dataset:"}\NormalTok{, df\_sample.shape)}
\BuiltInTok{print}\NormalTok{(}\StringTok{"}\CharTok{\textbackslash{}n}\StringTok{Column types:"}\NormalTok{)}
\BuiltInTok{print}\NormalTok{(df\_sample.dtypes.value\_counts())}
\BuiltInTok{print}\NormalTok{(}\StringTok{"}\CharTok{\textbackslash{}n}\StringTok{Missing values per column:"}\NormalTok{)}
\BuiltInTok{print}\NormalTok{(df\_sample.isna().}\BuiltInTok{sum}\NormalTok{().sort\_values(ascending}\OperatorTok{=}\VariableTok{False}\NormalTok{).head(}\DecValTok{10}\NormalTok{))}

\CommentTok{\# ============================}
\CommentTok{\# 4️⃣ Summary statistics for numeric variables}
\CommentTok{\# ============================}
\BuiltInTok{print}\NormalTok{(}\StringTok{"}\CharTok{\textbackslash{}n}\StringTok{📈 Summary statistics for numeric variables:"}\NormalTok{)}
\NormalTok{display(df\_sample.describe().T)}

\CommentTok{\# ============================}
\CommentTok{\# 5️⃣ Overview of categorical variables}
\CommentTok{\# ============================}
\BuiltInTok{print}\NormalTok{(}\StringTok{"}\CharTok{\textbackslash{}n}\StringTok{🔠 Categorical variable overview (top 10):"}\NormalTok{)}
\NormalTok{cat\_cols }\OperatorTok{=}\NormalTok{ df\_sample.select\_dtypes(include}\OperatorTok{=}\StringTok{"object"}\NormalTok{).columns}

\ControlFlowTok{for}\NormalTok{ col }\KeywordTok{in}\NormalTok{ cat\_cols[:}\DecValTok{10}\NormalTok{]:}
    \BuiltInTok{print}\NormalTok{(}\SpecialStringTok{f"}\CharTok{\textbackslash{}n}\SpecialStringTok{Column: }\SpecialCharTok{\{}\NormalTok{col}\SpecialCharTok{\}}\SpecialStringTok{"}\NormalTok{)}
    \BuiltInTok{print}\NormalTok{(}\StringTok{"Unique values:"}\NormalTok{, df\_sample[col].nunique())}
    \CommentTok{\# truncate long text values for readability}
    \BuiltInTok{print}\NormalTok{(df\_sample[col].astype(}\BuiltInTok{str}\NormalTok{).}\BuiltInTok{str}\NormalTok{.}\BuiltInTok{slice}\NormalTok{(}\DecValTok{0}\NormalTok{, }\DecValTok{100}\NormalTok{).value\_counts().head(}\DecValTok{3}\NormalTok{))}

\CommentTok{\# ============================}
\CommentTok{\# 6️⃣ Summary of cleaned fields}
\CommentTok{\# ============================}
\ControlFlowTok{if} \StringTok{"reaction\_clean"} \KeywordTok{in}\NormalTok{ df\_sample.columns }\KeywordTok{or} \StringTok{"drug\_clean"} \KeywordTok{in}\NormalTok{ df\_sample.columns:}
    \BuiltInTok{print}\NormalTok{(}\StringTok{"}\CharTok{\textbackslash{}n}\StringTok{🧪 Cleaned columns summary:"}\NormalTok{)}
    \ControlFlowTok{for}\NormalTok{ col }\KeywordTok{in}\NormalTok{ [}\StringTok{"drug\_clean"}\NormalTok{, }\StringTok{"reaction\_clean"}\NormalTok{]:}
        \ControlFlowTok{if}\NormalTok{ col }\KeywordTok{in}\NormalTok{ df\_sample.columns:}
            \BuiltInTok{print}\NormalTok{(}\SpecialStringTok{f"}\CharTok{\textbackslash{}n}\SpecialCharTok{\{}\NormalTok{col}\SpecialCharTok{\}}\SpecialStringTok{: }\SpecialCharTok{\{}\NormalTok{df\_sample[col]}\SpecialCharTok{.}\NormalTok{nunique()}\SpecialCharTok{\}}\SpecialStringTok{ unique values"}\NormalTok{)}
            \BuiltInTok{print}\NormalTok{(df\_sample[col].value\_counts(dropna}\OperatorTok{=}\VariableTok{True}\NormalTok{).head(}\DecValTok{5}\NormalTok{))}
\end{Highlighting}
\end{Shaded}

\begin{verbatim}
📊 Dataset Overview
========================================
Shape of dataset: (2000, 43)

Column types:
object     17
int64      13
float64    13
Name: count, dtype: int64

Missing values per column:
reportduplicate                             1995
seriousnesscongenitalanomali                1994
primarysource.literaturereference           1983
authoritynumb                               1958
seriousnessdisabling                        1911
patient.patientagegroup                     1909
seriousnesslifethreatening                  1900
seriousnessdeath                            1846
patient.summary.narrativeincludeclinical    1459
patient.patientweight                       1136
dtype: int64

📈 Summary statistics for numeric variables:
\end{verbatim}

\begin{longtable}[]{@{}lllllllll@{}}
\toprule\noalign{}
& count & mean & std & min & 25\% & 50\% & 75\% & max \\
\midrule\noalign{}
\endhead
\bottomrule\noalign{}
\endlastfoot
safetyreportversion & 2000.0 & 1.841000e+00 & 1.655926 & 1.00 &
1.000000e+00 & 2.0 & 2.000000e+00 & 48.0 \\
safetyreportid & 2000.0 & 1.011005e+07 & 113592.713798 & 10003332.00 &
1.003541e+07 & 10070505.5 & 1.013941e+07 & 10533593.0 \\
transmissiondateformat & 2000.0 & 1.020000e+02 & 0.000000 & 102.00 &
1.020000e+02 & 102.0 & 1.020000e+02 & 102.0 \\
transmissiondate & 2000.0 & 2.014546e+07 & 9528.358215 & 20141002.00 &
2.014100e+07 & 20141212.0 & 2.015033e+07 & 20240716.0 \\
reporttype & 2000.0 & 1.377500e+00 & 0.696589 & 1.00 & 1.000000e+00 &
1.0 & 2.000000e+00 & 4.0 \\
serious & 2000.0 & 1.204500e+00 & 0.403437 & 1.00 & 1.000000e+00 & 1.0 &
1.000000e+00 & 2.0 \\
seriousnesshospitalization & 901.0 & 1.002220e+00 & 0.047088 & 1.00 &
1.000000e+00 & 1.0 & 1.000000e+00 & 2.0 \\
receivedateformat & 2000.0 & 1.020000e+02 & 0.000000 & 102.00 &
1.020000e+02 & 102.0 & 1.020000e+02 & 102.0 \\
receivedate & 2000.0 & 2.014041e+07 & 591.668543 & 20120703.00 &
2.014032e+07 & 20140409.0 & 2.014043e+07 & 20141022.0 \\
receiptdateformat & 2000.0 & 1.020000e+02 & 0.000000 & 102.00 &
1.020000e+02 & 102.0 & 1.020000e+02 & 102.0 \\
receiptdate & 2000.0 & 2.014301e+07 & 8854.602633 & 20130524.00 &
2.014033e+07 & 20140502.0 & 2.014082e+07 & 20240527.0 \\
fulfillexpeditecriteria & 2000.0 & 1.302500e+00 & 0.459455 & 1.00 &
1.000000e+00 & 1.0 & 2.000000e+00 & 2.0 \\
duplicate & 1916.0 & 1.000000e+00 & 0.000000 & 1.00 & 1.000000e+00 & 1.0
& 1.000000e+00 & 1.0 \\
primarysource.qualification & 1900.0 & 3.014211e+00 & 1.743810 & 1.00 &
1.000000e+00 & 3.0 & 5.000000e+00 & 5.0 \\
sender.sendertype & 2000.0 & 2.000000e+00 & 0.000000 & 2.00 &
2.000000e+00 & 2.0 & 2.000000e+00 & 2.0 \\
receiver.receivertype & 2000.0 & 6.000000e+00 & 0.000000 & 6.00 &
6.000000e+00 & 6.0 & 6.000000e+00 & 6.0 \\
patient.patientonsetage & 1725.0 & 1.373194e+02 & 1387.931646 & 0.00 &
5.300000e+01 & 64.0 & 7.300000e+01 & 32402.0 \\
patient.patientonsetageunit & 1725.0 & 8.010191e+02 & 0.207695 & 801.00
& 8.010000e+02 & 801.0 & 8.010000e+02 & 804.0 \\
patient.patientsex & 1979.0 & 1.535624e+00 & 0.525506 & 0.00 &
1.000000e+00 & 2.0 & 2.000000e+00 & 2.0 \\
seriousnessother & 997.0 & 1.001003e+00 & 0.031670 & 1.00 & 1.000000e+00
& 1.0 & 1.000000e+00 & 2.0 \\
seriousnessdeath & 154.0 & 1.025974e+00 & 0.159577 & 1.00 & 1.000000e+00
& 1.0 & 1.000000e+00 & 2.0 \\
patient.patientweight & 864.0 & 8.185603e+01 & 23.434724 & 10.97 &
6.582250e+01 & 78.9 & 9.317500e+01 & 242.0 \\
seriousnesslifethreatening & 100.0 & 1.020000e+00 & 0.140705 & 1.00 &
1.000000e+00 & 1.0 & 1.000000e+00 & 2.0 \\
seriousnessdisabling & 89.0 & 1.056180e+00 & 0.231573 & 1.00 &
1.000000e+00 & 1.0 & 1.000000e+00 & 2.0 \\
seriousnesscongenitalanomali & 6.0 & 1.833333e+00 & 0.408248 & 1.00 &
2.000000e+00 & 2.0 & 2.000000e+00 & 2.0 \\
patient.patientagegroup & 91.0 & 5.373626e+00 & 0.550724 & 3.00 &
5.000000e+00 & 5.0 & 6.000000e+00 & 6.0 \\
\end{longtable}

\begin{verbatim}

🔠 Categorical variable overview (top 10):

Column: primarysourcecountry
Unique values: 47
primarysourcecountry
US    1320
GB     215
BR      90
Name: count, dtype: int64

Column: occurcountry
Unique values: 43
occurcountry
US     1103
nan     310
GB      193
Name: count, dtype: int64

Column: companynumb
Unique values: 1895
companynumb
nan                          83
US-PFIZER INC-2014107015      2
US-MYLANLABS-2014S1010381     2
Name: count, dtype: int64

Column: reportduplicate.duplicatesource
Unique values: 124
reportduplicate.duplicatesource
GLAXOSMITHKLINE    272
PFIZER             203
MERCK               94
Name: count, dtype: int64

Column: reportduplicate.duplicatenumb
Unique values: 1889
reportduplicate.duplicatenumb
nan                          89
163-21880-12122701            2
US-MYLANLABS-2014S1010381     2
Name: count, dtype: int64

Column: primarysource.reportercountry
Unique values: 46
primarysource.reportercountry
US                       1055
COUNTRY NOT SPECIFIED     316
GB                        215
Name: count, dtype: int64

Column: sender.senderorganization
Unique values: 1
sender.senderorganization
FDA-Public Use    2000
Name: count, dtype: int64

Column: receiver.receiverorganization
Unique values: 1
receiver.receiverorganization
FDA    2000
Name: count, dtype: int64

Column: patient.reaction
Unique values: 1744
patient.reaction
[{'reactionmeddraversionpt': '18.0', 'reactionmeddrapt': 'Myocardial infarction', 'reactionoutcome':    57
[{'reactionmeddraversionpt': '17.0', 'reactionmeddrapt': 'Drug ineffective', 'reactionoutcome': '6'}    29
[{'reactionmeddraversionpt': '18.0', 'reactionmeddrapt': 'Cerebrovascular accident', 'reactionoutcom    27
Name: count, dtype: int64

Column: patient.drug
Unique values: 1976
patient.drug
[{'drugcharacterization': '1', 'medicinalproduct': 'AVANDIA', 'drugauthorizationnumb': '021071', 'dr    72
[{'drugcharacterization': '1', 'medicinalproduct': 'PRADAXA', 'drugauthorizationnumb': '022512', 'dr    27
[{'drugcharacterization': '1', 'medicinalproduct': 'LYRICA', 'drugauthorizationnumb': '021446', 'dru    22
Name: count, dtype: int64

🧪 Cleaned columns summary:

drug_clean: 820 unique values
drug_clean
AVANDIA        77
JAKAFI         57
ASPIRIN.       40
PROPRANOLOL    39
REVLIMID       36
Name: count, dtype: int64

reaction_clean: 751 unique values
reaction_clean
Myocardial infarction       102
Cerebrovascular accident     50
Drug ineffective             41
Death                        26
Fatigue                      20
Name: count, dtype: int64
\end{verbatim}

\subsection{Questions for reviewers}\label{questions-for-reviewers}

\begin{enumerate}
\def\labelenumi{\arabic{enumi}.}
\tightlist
\item
  Is the project question, predicting the type of adverse reaction for
  cardiovascular drugs using FAERS data, clearly defined and feasible
  for a classical ML model?
\item
  Does using reaction type as the outcome variable make the analysis
  interpretable and meaningful for this dataset?
\item
  Do you think the sample size (a reduced FAERS dataset under 100 MB, I
  had to take a random subet of my OG data because it was far over 100MB
  ) is large enough to yield reliable results while still being
  computationally feasible?
\end{enumerate}

\section{Data 2}\label{data-2}

\subsection{Problem or question}\label{problem-or-question-1}

\begin{itemize}
\tightlist
\item
  Identify the problem you will solve or the question you will answer
\item
  Explain why you think this topic is important.
\item
  Identify the types of data/variables you will use.
\item
  State the major deliverable(s) you will create to solve this
  problem/answer this question.
\end{itemize}

\subsection{Introduction and data}\label{introduction-and-data-1}

If you are using a dataset:

\begin{itemize}
\item
  Identify the source of the data.
\item
  State when and how it was originally collected (by the original data
  curator, not necessarily how you found the data).
\item
  Write a brief description of the observations.
\item
  Address ethical concerns about the data, if any.
\end{itemize}

\subsection{Glimpse of data}\label{glimpse-of-data-1}

\subsection{Questions for reviewers}\label{questions-for-reviewers-1}

List specific questions for your peer reviewers and project mentor to
answer in giving you feedback on this topic.

\section{Data 3}\label{data-3}

\subsection{Problem or question}\label{problem-or-question-2}

\textbf{Problem:} Can simple self-reported survey answers (no labs)
identify adults who are likely to meet clinical prediabetes thresholds,
so they can be given secondary testing?

\textbf{Objective:} Train a calibrated classifier that estimates
P(prediabetes) using only questionnaire items available and use it to
determine who to give blood tests to.

\textbf{Why this matters:} Prediabetes is common and underdiagnosed. A
short survey would be low-cost and help public-health programs and
community clinics determine who should receive lab testing.

\textbf{Variable Types:}

\begin{itemize}
\item
  Quantitative: age, self-reported BMI, hours of sleep, weekly physical
  activity minutes.
\item
  Categorical/ordinal: sex, race/ethnicity, education, income-to-poverty
  ratio band, health insurance, family history of diabetes, self-rated
  health, smoking status, alcohol use, leisure/transport
  physical-activity frequency, diet questions (e.g., sugary beverage
  frequency).
\item
  Target from blood samples: Prediabetes = 1 if either A1c 5.7--6.4\%,
  or Fasting plasma glucose 100--125 mg/dL,
\item
  Usability Plan: We aim to build a small web app (Quarto + shinylive)
  where a user answers 10--15 survey questions and gets their predicted
  risk of being prediabetic along with a short explanation of
  contributing factors, and links suggesting follow up testing to
  confirm a diagnosis if their probability is high.
\end{itemize}

\subsection{Introduction and data}\label{introduction-and-data-2}

\textbf{Source:} NHANES (National Health and Nutrition Examination
Survey), run by NCHS/CDC.

\textbf{Collection:} Continuous, national survey since 1999 using
multistage, stratified sampling. The data comes from an in-home
questionnaire and a mobile examination center that collects lab samples.
Each 2-year cycle has its own design variables and weights. Years used:
2015--2016 and 2017--2018. We'll stack cycles, use variables that are
constant, and adjust weights per NHANES documentation.

\textbf{Observations:}

\begin{verbatim}
Dataset includes weights that estimate the amount of people a unique respondent represents 

Missingness varies by feature with some entirely missing (probably a mistake on my part)

Fasting glucose is less complete than A1c

Dataset reflects participant at a single time point 

The target distribution is skewed towards no diabetes.
\end{verbatim}

\textbf{Ethical Considerations:} The main ethical concern I can see that
may arise from the data would be if the tool were used for diagnostics
in a way that would replace lab testing or doctor visits. To address
this, we will put a warning that the tool is not diagnostic and is not
meant to be taken as medical advice. The other concern would be the
privacy of survey respondents. To address this, we will not save any
identifying information.

\subsection{Glimpse of data}\label{glimpse-of-data-2}

\begin{Shaded}
\begin{Highlighting}[]
\ImportTok{import}\NormalTok{ pandas }\ImportTok{as}\NormalTok{ pd}
\ImportTok{import}\NormalTok{ numpy }\ImportTok{as}\NormalTok{ np}

\NormalTok{target }\OperatorTok{=} \StringTok{"prediabetes"}
\NormalTok{df }\OperatorTok{=}\NormalTok{ pd.read\_csv(}\StringTok{"data{-}raw/nhanes1518\_selected\_labels.csv"}\NormalTok{)}
\NormalTok{df }\OperatorTok{=}\NormalTok{ df.drop(}
\NormalTok{    columns}\OperatorTok{=}\NormalTok{[}
        \StringTok{"RIDRETH3\_code"}\NormalTok{,}
        \StringTok{"RIAGENDR\_code"}\NormalTok{,}
        \StringTok{"RIDRETH1\_code"}\NormalTok{,}
        \StringTok{"DMDEDUC\_code"}\NormalTok{,}
        \StringTok{"INSURE\_code"}\NormalTok{,}
\NormalTok{    ],}
\NormalTok{    errors}\OperatorTok{=}\StringTok{"ignore"}\NormalTok{,}
\NormalTok{)}
\BuiltInTok{print}\NormalTok{(}\StringTok{"Shape"}\NormalTok{)}
\BuiltInTok{print}\NormalTok{(}\SpecialStringTok{f"Rows: }\SpecialCharTok{\{}\NormalTok{df}\SpecialCharTok{.}\NormalTok{shape[}\DecValTok{0}\NormalTok{]}\SpecialCharTok{:,\}}\SpecialStringTok{ | Columns: }\SpecialCharTok{\{}\NormalTok{df}\SpecialCharTok{.}\NormalTok{shape[}\DecValTok{1}\NormalTok{]}\SpecialCharTok{:,\}}\SpecialStringTok{"}\NormalTok{)}


\CommentTok{\# missingness}
\NormalTok{miss }\OperatorTok{=}\NormalTok{ (}
\NormalTok{    pd.DataFrame(}
\NormalTok{        \{}
            \StringTok{"column"}\NormalTok{: df.columns,}
            \StringTok{"missing"}\NormalTok{: df.isna().}\BuiltInTok{sum}\NormalTok{().values,}
            \StringTok{"missing\_pct"}\NormalTok{: (df.isna().mean().values }\OperatorTok{*} \DecValTok{100}\NormalTok{).}\BuiltInTok{round}\NormalTok{(}\DecValTok{2}\NormalTok{),}
            \StringTok{"dtype"}\NormalTok{: [df[c].dtype }\ControlFlowTok{for}\NormalTok{ c }\KeywordTok{in}\NormalTok{ df.columns],}
\NormalTok{        \}}
\NormalTok{    )}
\NormalTok{    .sort\_values([}\StringTok{"missing\_pct"}\NormalTok{, }\StringTok{"column"}\NormalTok{], ascending}\OperatorTok{=}\NormalTok{[}\VariableTok{False}\NormalTok{, }\VariableTok{True}\NormalTok{])}
\NormalTok{    .reset\_index(drop}\OperatorTok{=}\VariableTok{True}\NormalTok{)}
\NormalTok{)}
\NormalTok{display(miss)}

\CommentTok{\# numeric}
\NormalTok{num\_cols }\OperatorTok{=}\NormalTok{ df.select\_dtypes(include}\OperatorTok{=}\NormalTok{[np.number, }\StringTok{"bool"}\NormalTok{]).columns.tolist()}
\ControlFlowTok{if}\NormalTok{ num\_cols:}
\NormalTok{    num\_desc }\OperatorTok{=}\NormalTok{ df[num\_cols].describe().T}
\NormalTok{    num\_desc[}\StringTok{"missing\_pct"}\NormalTok{] }\OperatorTok{=}\NormalTok{ df[num\_cols].isna().mean().values }\OperatorTok{*} \DecValTok{100}
\NormalTok{    display(num\_desc.}\BuiltInTok{round}\NormalTok{(}\DecValTok{1}\NormalTok{))}

    \CommentTok{\# Target Distribution}
\ControlFlowTok{if}\NormalTok{ target }\KeywordTok{is} \VariableTok{None}\NormalTok{:}
    \ControlFlowTok{for}\NormalTok{ cand }\KeywordTok{in}\NormalTok{ [}
        \StringTok{"prediabetes"}\NormalTok{,}
        \StringTok{"prediabetes\_label"}\NormalTok{,}
        \StringTok{"preDM"}\NormalTok{,}
        \StringTok{"pre\_dm"}\NormalTok{,}
        \StringTok{"pre\_diabetes"}\NormalTok{,}
        \StringTok{"label"}\NormalTok{,}
        \StringTok{"y"}\NormalTok{,}
\NormalTok{    ]:}
        \ControlFlowTok{if}\NormalTok{ cand }\KeywordTok{in}\NormalTok{ df.columns:}
\NormalTok{            target }\OperatorTok{=}\NormalTok{ cand}
            \ControlFlowTok{break}
\ControlFlowTok{if}\NormalTok{ target }\KeywordTok{and}\NormalTok{ target }\KeywordTok{in}\NormalTok{ df.columns:}
    \BuiltInTok{print}\NormalTok{(}\SpecialStringTok{f"}\CharTok{\textbackslash{}n}\SpecialStringTok{Target dirstribution: }\SpecialCharTok{\{}\NormalTok{target}\SpecialCharTok{\}}\SpecialStringTok{ "}\NormalTok{)}
\NormalTok{    t }\OperatorTok{=}\NormalTok{ df[target]}
    \ControlFlowTok{if}\NormalTok{ t.dropna().dtype }\OperatorTok{==} \BuiltInTok{bool}\NormalTok{:}
\NormalTok{        t\_bin }\OperatorTok{=}\NormalTok{ t.astype(}\BuiltInTok{int}\NormalTok{)}
    \ControlFlowTok{elif}\NormalTok{ t.dtype.kind }\KeywordTok{in} \StringTok{"biufc"}\NormalTok{:}
\NormalTok{        t\_bin }\OperatorTok{=}\NormalTok{ t}
    \ControlFlowTok{else}\NormalTok{:}
\NormalTok{        mapping }\OperatorTok{=}\NormalTok{ \{}\StringTok{"yes"}\NormalTok{: }\DecValTok{1}\NormalTok{, }\StringTok{"no"}\NormalTok{: }\DecValTok{0}\NormalTok{, }\StringTok{"true"}\NormalTok{: }\DecValTok{1}\NormalTok{, }\StringTok{"false"}\NormalTok{: }\DecValTok{0}\NormalTok{, }\StringTok{"y"}\NormalTok{: }\DecValTok{1}\NormalTok{, }\StringTok{"n"}\NormalTok{: }\DecValTok{0}\NormalTok{\}}
\NormalTok{        t\_bin }\OperatorTok{=}\NormalTok{ t.astype(}\BuiltInTok{str}\NormalTok{).}\BuiltInTok{str}\NormalTok{.lower().}\BuiltInTok{map}\NormalTok{(mapping)}

\NormalTok{    counts }\OperatorTok{=}\NormalTok{ t\_bin.value\_counts(dropna}\OperatorTok{=}\VariableTok{False}\NormalTok{).to\_frame(name}\OperatorTok{=}\StringTok{"count"}\NormalTok{)}
\NormalTok{    counts[}\StringTok{"pct"}\NormalTok{] }\OperatorTok{=}\NormalTok{ (counts[}\StringTok{"count"}\NormalTok{] }\OperatorTok{/} \BuiltInTok{len}\NormalTok{(df) }\OperatorTok{*} \DecValTok{100}\NormalTok{).}\BuiltInTok{round}\NormalTok{(}\DecValTok{2}\NormalTok{)}
\NormalTok{    display(counts)}
\NormalTok{    prev }\OperatorTok{=} \BuiltInTok{float}\NormalTok{(np.nanmean(pd.to\_numeric(t\_bin, errors}\OperatorTok{=}\StringTok{"coerce"}\NormalTok{)))}
    \BuiltInTok{print}\NormalTok{(}\SpecialStringTok{f"Prevalence (mean of 0/1) ≈ }\SpecialCharTok{\{}\NormalTok{prev}\SpecialCharTok{:.4f\}}\SpecialStringTok{"}\NormalTok{)}
\end{Highlighting}
\end{Shaded}

\begin{verbatim}
Shape
Rows: 8,880 | Columns: 20
\end{verbatim}

\begin{longtable}[]{@{}lllll@{}}
\toprule\noalign{}
& column & missing & missing\_pct & dtype \\
\midrule\noalign{}
\endhead
\bottomrule\noalign{}
\endlastfoot
0 & sleep\_hours & 8880 & 100.00 & float64 \\
1 & alc\_unit\_ans & 6098 & 68.67 & object \\
2 & alc\_qty & 5507 & 62.02 & float64 \\
3 & smoke\_now\_ans & 5391 & 60.71 & object \\
4 & alc\_days\_yr & 4756 & 53.56 & float64 \\
5 & pir & 1024 & 11.53 & float64 \\
6 & educ & 480 & 5.41 & object \\
7 & SEQN & 0 & 0.00 & float64 \\
8 & age & 0 & 0.00 & int64 \\
9 & cycle & 0 & 0.00 & object \\
10 & fam\_hx\_dm\_ans & 0 & 0.00 & object \\
11 & insured & 0 & 0.00 & object \\
12 & physact\_min & 0 & 0.00 & float64 \\
13 & prediabetes & 0 & 0.00 & int64 \\
14 & race\_eth & 0 & 0.00 & object \\
15 & sex & 0 & 0.00 & object \\
16 & sleep\_trouble\_ans & 0 & 0.00 & object \\
17 & smoke\_ever\_ans & 0 & 0.00 & object \\
18 & soda\_freq & 0 & 0.00 & float64 \\
19 & weight & 0 & 0.00 & float64 \\
\end{longtable}

\begin{longtable}[]{@{}llllllllll@{}}
\toprule\noalign{}
& count & mean & std & min & 25\% & 50\% & 75\% & max & missing\_pct \\
\midrule\noalign{}
\endhead
\bottomrule\noalign{}
\endlastfoot
SEQN & 8880.0 & 93383.4 & 5528.1 & 83733.0 & 88607.2 & 93482.5 & 98126.2
& 102956.0 & 0.0 \\
prediabetes & 8880.0 & 0.4 & 0.5 & 0.0 & 0.0 & 0.0 & 1.0 & 1.0 & 0.0 \\
weight & 8880.0 & 22849.6 & 23700.7 & 2289.8 & 9069.6 & 13715.4 &
25432.2 & 209881.4 & 0.0 \\
age & 8880.0 & 46.8 & 18.3 & 18.0 & 31.0 & 46.0 & 61.0 & 80.0 & 0.0 \\
pir & 7856.0 & 2.5 & 1.6 & 0.0 & 1.1 & 2.1 & 4.0 & 5.0 & 11.5 \\
alc\_days\_yr & 4124.0 & 1.3 & 0.5 & 1.0 & 1.0 & 1.0 & 2.0 & 9.0 &
53.6 \\
alc\_qty & 3373.0 & 3.8 & 26.7 & 0.0 & 1.0 & 2.0 & 4.0 & 999.0 & 62.0 \\
physact\_min & 8880.0 & 5.7 & 3.1 & 4.0 & 4.0 & 4.0 & 7.0 & 106.0 &
0.0 \\
sleep\_hours & 0.0 & NaN & NaN & NaN & NaN & NaN & NaN & NaN & 100.0 \\
soda\_freq & 8880.0 & 3.0 & 1.0 & 1.0 & 2.0 & 3.0 & 4.0 & 9.0 & 0.0 \\
\end{longtable}

\begin{verbatim}

Target dirstribution: prediabetes 
\end{verbatim}

\begin{longtable}[]{@{}lll@{}}
\toprule\noalign{}
& count & pct \\
prediabetes & & \\
\midrule\noalign{}
\endhead
\bottomrule\noalign{}
\endlastfoot
0 & 4996 & 56.26 \\
1 & 3884 & 43.74 \\
\end{longtable}

\begin{verbatim}
Prevalence (mean of 0/1) ≈ 0.4374
\end{verbatim}

\subsection{Questions for reviewers}\label{questions-for-reviewers-2}

List specific questions for your peer reviewers and project mentor to
answer in giving you feedback on this topic.

\begin{itemize}
\item
  I had some trouble collecting data from the website and putting it
  together into a csv. Heres the website:
  https://wwwn.cdc.gov/nchs/nhanes/search/datapage.aspx?Component=Questionnaire\&CycleBeginYear=2015
  Do you have any reccomendations for resources that may help in this
  process?
\item
  I would also appreciate guidance on how to work with the sampling
  weights that tell us how many people in the U.S. each respondent
  represents.
\end{itemize}




\end{document}
